\documentclass[12pt]{article}
\usepackage[utf8]{inputenc}
\usepackage[T1]{fontenc}
\usepackage[margin=1in]{geometry}
\usepackage{setspace}
\usepackage{xcolor}
\usepackage[colorlinks=true, urlcolor=blue!60!black, linkcolor=blue!60!black]{hyperref}
\setlength{\parindent}{0em}
\setlength{\parskip}{0.45em}
\renewcommand{\baselinestretch}{1.0}
\begin{document}
\begin{center}
{\Large \textit{Texas v. Johnson}, 491 U.S. 397 (1989)\par}
\vspace{0.3em}
{\normalsize Case Brief --- One-Page Edition\par}
\end{center}

\textbf{Memory Jogger.} The flag-burning case: offensive political expression may not be punished for its message (\href{https://www.courtlistener.com/opinion/112304/texas-v-johnson/}{491 U.S. 397}).

\textbf{Facts.} At the 1984 Republican National Convention, Gregory Lee Johnson burned an American flag outside Dallas City Hall in protest. No one was hurt, but onlookers were offended, and Texas convicted him of flag desecration.

\textbf{Procedural History.} The Texas Court of Criminal Appeals reversed his conviction; the Supreme Court granted certiorari.

\textbf{Votes.} Affirmed, 5--4. Brennan wrote for the Court; Kennedy concurred; Rehnquist, White, O'Connor, and Stevens dissented.

\textbf{Holding.} Flag burning is expressive conduct, and because Texas's interest in the flag's symbolic meaning is tied to suppressing expression, the conviction fails the most exacting scrutiny.

\textbf{Reasoning.} No breach of the peace occurred, and offense cannot be presumed to incite violence. The law punished only offensive burning, and government may not prescribe orthodoxy for national symbols. Kennedy noted the ``personal toll'' of a compelled result; Rehnquist called flag burning ``an inarticulate grunt or roar''; Stevens saw only a ``trivial burden'' on speakers.

\textbf{Examples.} \emph{Same-Side:} the federal Flag Protection Act (\textit{Eichman}). \emph{Opposite-Side:} prosecuting theft of a flag or a fire-code violation. \emph{Fence-Sitter:} flag burning amid threats at a military funeral.

\textbf{Critique.} Praised for content neutrality; criticized for undervaluing civic cohesion and the original meaning of ``speech.''

\textbf{Key Quotation.} ``If there is a bedrock principle underlying the First Amendment, it is that the government may not prohibit the expression of an idea simply because society finds the idea itself offensive or disagreeable.'' --- Justice Brennan, Opinion of the Court, 491 U.S. at 414.

\end{document}
